\documentclass[../../main.tex]{subfiles}
\begin{document}
\begin{oframed}
  \subfile{problem}
\end{oframed}

{\bf [解答]}

\begin{figure}[htb]
  \centering
  \begin{tikzpicture}[x=0.9cm, y=0.55cm, >=stealth, every node/.style={font=\small}]
    % 座標軸
    \draw[->, thick, gray!80!black] (-7.2, 0) -- (3.8, 0) node[right] {$x$};
    \draw[->, thick, gray!80!black] (0, -0.6) -- (0, 10.0) node[above] {$y$};
    \node[below right=2pt] at (0, 0) {$O$};

    % 放物線 C_1: y = x^2
    \draw[line width=1.2pt, blue!85!black, domain=-3.1:3.1, samples=60]
    plot (\x, {(\x)^2}) node[above right] {$C_1: y=x^2$};

    % 放物線 C_2: y = 3(x+4)^2 + 2
    \draw[line width=1.2pt, orange!85!black, domain=-5.6:-2.4, samples=60]
    plot (\x, {3*(\x + 4)^2 + 2}) node[above left] {$C_2: y=3x^2+24x+50$};

    % 点 P(-6, 3)
    \coordinate (P) at (-6, 3);
    \node[circle, fill=red!85!black, inner sep=1.4pt] at (P) {};
    \node[above left=1pt, red!85!black] at (P) {$P(-6,3)$};

    % 相似の中心からの直線 1: P -> B_0(-4, 2) -> A_0(0, 0)
    \coordinate (B0) at (-4, 2);
    \coordinate (A0) at (0, 0);
    \draw[line width=1.0pt, dashed, gray!70!black] (P) -- (A0);
    \node[circle, fill=orange!85!black, inner sep=1.2pt] at (B0) {};
    \node[below=2pt, font=\footnotesize] at (B0) {$B_0(-4,2)$};

    % 相似の中心からの直線 2: P -> B_1(-3, 5) -> A_1(3, 9)
    \coordinate (B1) at (-3, 5);
    \coordinate (A1) at (3, 9);
    \draw[line width=1.0pt, dashed, gray!70!black] (P) -- (A1);
    \node[circle, fill=orange!85!black, inner sep=1.2pt] at (B1) {};
    \node[left=2pt, font=\footnotesize] at (B1) {$B$};
    \node[circle, fill=blue!85!black, inner sep=1.2pt] at (A1) {};
    \node[above right=1pt, font=\footnotesize] at (A1) {$A(3,9)$};

    % 内分比の注記
    \node[above=1pt, font=\scriptsize, gray!80!black] at ($ (A1)!0.33!(B1) $) {$m=2$};
    \node[above=1pt, font=\scriptsize, gray!80!black] at ($ (B1)!0.5!(P) $) {$n=1$};
  \end{tikzpicture}
  \caption{放物線 $C_1, C_2$ と相似の中心 $P(-6,3)$ による内分関係}
  \label{fig:homothety_parabolas}
\end{figure}

題意より$t\in\mathbb{R}$に対して$A(t,t^2)$と置ける． 題意の条件から$s\in\mathbb{R}$に対して$B(s,3s^2+24s+50)$とおいてよい．このとき内分点に関する条件から
\begin{numcases}{}
  \dfrac{nt+m\alpha}{m+n}=s \label{3}\\
  \dfrac{nt^2+m\beta}{m+n}=3s^2+24s+50 \label{4}
\end{numcases}
$\forall t \exists s,\cref{3}\land\cref{4}$となる$P$，$m$，$n$の条件を求めればよい．
\cref{3}を\cref{4}に代入して$s$を消去する．$a=m+n$として簡単のため
\begin{align}
  p&=\dfrac{6n(m\alpha+4a)}{a^2} \\
  q&=\dfrac{3m^2\alpha^2}{a^2}+\dfrac{24m\alpha}{a}+50
\end{align}
とおけば，
\begin{align}
  &\dfrac{nt^2+m\beta}{a}\\
  &=3\left(\dfrac{nt+m\alpha}{a}\right)^2+24 \dfrac{nt+m\alpha}{a}+50  \\
  \Longleftrightarrow&\dfrac{n}{a}t^2+\dfrac{m\beta}{a}=\dfrac{3n^2}{a^2}t^2+pt+q
\end{align}
が恒等式であればよい． したがって
\begin{align}
  \begin{cases}
    \dfrac{n}{a}=\dfrac{3n^2}{a^2}  \\
    p=0  \\
    \dfrac{m\beta}{a}=q
  \end{cases}
  \Longleftrightarrow
  \begin{cases}
    1=\dfrac{3n}{a}  \\
    m\alpha+4a=0  \\
    \dfrac{m\beta}{a}=q
  \end{cases} \quad (\because a,n\neq 0)
\end{align}
となる．第$1$式および$a=m+n$から，$n:m:a=1:2:3$である．これを第$2$式に代入して$\alpha=-6$を得る．
さらに第$3$式にこれらを代入して
\begin{align}
  &\dfrac{2}{3}\beta=\dfrac{4}{3}\alpha^2+16\alpha+50  \\
  \therefore &\beta=3
\end{align}
である．以上から$P(-6,3)$，$m:n=2:1$である．$\cdots$(答)
\newpage
\end{document}
